Die automatisierte Erkennung von Wildtierarten auf einem Handgerät muss Hunderte visuell ähnlicher Säugetierarten innerhalb eines festen Zeitbudgets pro Bild unterscheiden, und zwar auf Hardware, die für die in diesem Feld üblichen Ensembles aus Detektor und Klassifikator zu klein ist. Naheliegend ist ein Detektor im Nano-Maßstab, und naheliegend für dessen Training ist die Wissensdestillation aus einem größeren Lehrermodell. Diese Arbeit prüft diese Annahme. Sie untersucht, ob die Destillation eines feinabgestimmten Lehrermodells in ein Schülermodell im Nano-Maßstab eine höhere Genauigkeit erzielt als das direkte Feinabstimmen desselben Schülermodells auf der Zieldomäne, über ein Klassenuniversum von $225$ Säugetierarten, von denen $50$ jeweils weniger als $150$ verwendbare Fotografien aufweisen.

Offene Quellen lieferten einen Rohkorpus unter einer auf kommerzielle Nutzung ausgerichteten Lizenzpolitik, den eine mehrstufige Filterung aus Metadaten-Heuristiken, einem Kamerafallen-Detektor und einem Artklassifikator auf $222{.}073$ reale Bilder über die $225$ Klassen hinweg kuratierte. Klassen mit zu wenigen Bildern wurden durch generierte Aufnahmen ergänzt, geordnet nach einer Bandstruktur, die die Zusammensetzung der Trainingsdaten bei gleichem Bildbudget von vollständig synthetisch bis vollständig real variiert. Eine eigenständige Teilstudie verglich zwölf Kombinationen aus Generator und Prompt-Regime, die jeweils $1{.}200$ synthetische Bilder derselben $12$ Klassen an dasselbe Detektormodell lieferten. Bewertet wurde auf einem gemischten Testsatz aus realen und synthetischen Bildern, durchgängig berichtet gemeinsam mit einer Auswertung allein auf den realen Bildern. Die Latenz wurde auf einem \textit{Raspberry Pi 400} als Stellvertreter für den Zielchip \textit{Qualcomm QCS605} gemessen.

Die Destillation unterlag. Das direkte Feinabstimmen lag in allen $17$ berichteten Zellen vorn, mit $0{,}599$ mAP auf dem gemischten Testsatz gegenüber $0{,}560$ und $0{,}529$ gegenüber $0{,}479$ allein auf den realen Bildern, mit dem größten Abstand auf dem datenarmen Band, auf dem die weichen Zielverteilungen am ehesten hätten helfen sollen. Gemessen wurde eine einzige Einstellung aus Destillationstemperatur und Gewichtung der Verlustanteile, bei einem Parameterverhältnis von Lehrer zu Schüler von etwa $72$ zu $1$, sodass das Ergebnis eine konkrete Rezeptur begrenzt und nicht das Verfahren. Keines der beiden vor dem Training verfügbaren Kriterien, weder der scheinbare Realismus noch der Preis pro Bild, sagte vorher, welcher Generator die nützlichsten Daten liefert, während der Detailgrad des Prompts die Wahl des Generators übertraf und ein gekürzter Prompt etwa die Hälfte der Genauigkeit kostete. Die Hälfte eines kleinen realen Trainingssatzes durch synthetische Bilder zu ersetzen kostete nichts Messbares, mit $0{,}588$ gegenüber $0{,}591$ mAP, während der vollständige Ersatz etwa die Hälfte kostete. Eine vorab festgelegte Kontrollbedingung für den gemischten Testsatz wurde ausgelöst, die daraufhin vorgesehene Revision jedoch nicht mehr durchgeführt, die einzige hier offengelegte Abweichung vom Protokoll. Der resultierende Detektor benötigt weiterhin $423$\,ms pro Bild und damit das Zwölffache des Budgets von $30$\,ms, sodass nicht die Genauigkeit den Einsatz blockiert, sondern die Kompression.
